% Standalone with pdflatex. For inclusion in another document take the body.
\documentclass[11pt]{article}
\usepackage{amsmath,amsthm}
\newtheorem{theorem}{Theorem}
\newcommand{\wordalph}{\operatorname{alph}}
\newcommand{\Imp}{\operatorname{Imp}}
\newcommand{\Dist}{\mathcal D}
\begin{document}
\title{Short distribution representatives for locally counting targets}
\author{R14 research note}\date{23 November 2026}\maketitle

Put $q=\ell-1$, $k=|\Sigma|\ge1$, $V=k^q$, $K=k^\ell$.
Write $b(w)$ for the word tagged short if $|w|<\max(1,2q)$, otherwise
the tagged pair (first $q$, last $q$ letters). The type of a concatenation is
determined by the types of its pieces, including the empty word.
Let $T=|b(\Sigma^+)|$. Write $N(w)$ for the vector of multiplicities of
length-$\ell$ factors, counting overlaps, and $c(w)=\wordalph(w)$.
A typing is a homomorphism on positive words to a finite semigroup.
The order $s\preceq(c,b)$ means that equality on the right implies equality on the left.

An admissible context of arity $d$ inserts, in order, the $d$ nonempty
components of a tuple once each. Internal gaps are positive, the two
outer gaps arbitrary. Write $\Dist(x)$ for the accepting contexts of $x$.
Equivalence means equality both of component $s$-types and of distributions.
A target is $(f,s)$-substitutable if for every $1\le d\le f$,
two equally typed tuples sharing an accepting context are equivalent.

A grammar rule builds its output tuple by concatenating terminals and
child-component variables. Good rules have nonempty template components,
use each variable exactly once, use each child's variables in their indexed
order in the flattened output, and never place two consecutive variables of the
same child adjacent in one output component. The start symbol has one component.
Reduced means reachable and productive. Let $v$ count nonterminals, $r$ bound
the number of children, and $f$ bound the arity. Let $\theta$ be the maximum
of 1 and, over original rules, the minimum length of a successful sentence
using the rule. Each rule writes at most $\theta$ terminals, and each
nonterminal $B$ has a positive generated tuple $u_B$ of length at most
$\theta$: take a short successful witness for one of its rules.
A reachable occurrence of $B$ has a fixed accepting context $E_B$ for all
its generated tuples. Indeed successively filling siblings preserves hole
order, and a separator is either a terminal, a positive sibling piece, or
the inherited positive gap between two parent components. Consequently all
yields of a realized typed nonterminal $X$ (original symbol plus vector of
output types, in a successful derivation) belong to one equivalence class.
Let $\tau$ be the maximum, over these $X$, of the minimum total length
in that class, allowing representatives outside the grammar yield of $X$.

\begin{theorem}
Let $L$ be an $(f,s)$-substitutable $f$-MCFL. Suppose positive membership
depends only on $(b(w),N(w))$ and $s\preceq(c,b)$. For every reduced good
presentation of its positive part, with $f,r,\ell\ge1$, set
\[
 \lambda=K(2V-1),\quad B=2q\lambda,\quad c_*=c_{K+2},
\]
\[
 c_1=1,\quad c_{i+1}=Kc_i+f+B+2\quad(1\le i\le K+1).
\]
With $H=v((2^k-1)T+1)^f$, set
\[
 N_{\rm sk}=\ell c_*KfH,\qquad C_L=T^f(1+rH(\theta+f\ell))^{Kf}.
\]
Then
\[
 \tau\le f\ell+\ell N_{\rm sk}\theta+2\ell frN_{\rm sk}
             +\ell rN_{\rm sk} C_L H\theta
 =O_{k,\ell,f,r}(v^{Kf+2}\theta^{Kf+1}).
\]
\end{theorem}
\begin{proof}
\emph{Context invariant.}
A crossing window consists of a suffix of the first piece it touches,
a prefix of the last, and all intervening nonempty pieces. More precisely,
enumerate first and last nonempty piece and the two positive lengths taken,
each at most $q$; all intervening lengths and letters of the required
short pieces are determined by their tags. This counts occurrences with
multiplicities, also for a window spanning several boundaries. Thus equal
component $b,c$ and equal sum of internal $N$ imply equivalence in all
admissible outer contexts. The same boundary identity works for arbitrary
fixed pieces.

\emph{Tree and positions.}
Choose a derivation with fewest nodes for the output label $X$, call its
tuple $x$. An ancestor and proper descendant with equal original symbol
and output types could be contracted, preserving $X$ by the morphism property.
The height in nodes is at most $H$. Positions always refer to this original
tuple. By linear nondeletion every subtree component has a contiguous image
in one root component; incomparable subtrees have disjoint images.

A word of length at least $q$ is its initial $q$ letters followed by the
last-letter labels of a path in the de Bruijn graph with vertices $\Sigma^q$
and edges $\Sigma^\ell$ from prefix to suffix. Counts are edge counts. For each
root component select an empty one of the $K+1$ intervals $[c_g,c_{g+1})$
of frequencies. Call edges of frequency at least $c_{g+1}$ frequent,
the rest rare. If there is a frequent edge, there are at most
$R=(K-1)(c_g-1)$ rare occurrences. Every weak component $J$ of the
frequent edges is strongly connected: for an edge from which there is no
return, the vertices reachable from its head along frequent edges exclude
its tail. At least $c_{g+1}$ entries into this set require at least
$c_{g+1}-1>R$ exits in the full path, all rare. Every maximal interval
of frequent edges is in a single $J$. There are at most $R+1$ such
intervals for each $J$. Without frequent edges no pool is needed.

Mark every letter of the first $\min(c_*,N(x_i)(e))$ occurrences of every
edge $e$ in $x_i$. The skeleton consists of the writing nodes and their
ancestors, at most $N_{\rm sk}$ nodes. If empty, there are no windows
and $\|x\|_1\le fq$. Otherwise there are at most $rN_{\rm sk}$ maximal
off-skeleton subtrees. No internal window of one of their components
meets a mark and all such windows belong to one frequent interval.
A short-tag root component has no deletable window: all its windows, if
any, are selected ($q\ge1$ gives $c_*>2q$). Keep it unchanged on the
representative side. Components shorter than $\ell$ contribute no counts.

\emph{Cuts and colours.}
Fix one maximal off-skeleton subtree $\eta$ with components $z_\kappa$.
Order its leaves by depth first search with any fixed child order.
After the first $j$ leaves, $I_j$ consists of maximal subtrees all of whose
leaves are in that prefix (empty initially, $\{\eta\}$ at the end).
There are at most $rH$ of them: before completion they are completed
siblings on the next leaf's ancestor path. Their original terminal positions
form a nested sequence $U_j$. Newly covered positions are written on the
path of the newly completed leaf and number at most $H\theta$.
In the original derivation replace each member of $I_j$ by a derivation
of $u_B$ from its own original symbol. Write $z^j_\kappa$ for the new
components of $\eta$ and $W_j$ for the resulting sentence in $E_A$,
where $A$ underlies $X$. Each $W_j$ is generated, without any type
assumption on the replacements.

Tag every inserted letter. In each $z^j_\kappa$ record $b(z^j_\kappa)$
and the vector $\Imp_\kappa(j)$ of internal windows touching an inserted
letter. All pieces replacing original components are positive. Hence a
window with no inserted letter lies in one contiguous surviving stretch
and corresponds bijectively, by positions, to an original internal window
avoiding $U_j$. Conversely every such original window survives. A replacement
of length $n$ accounts for at most $n+q$ impure windows. Across the at most
$rH$ nodes the total is at most $rH(\theta+fq)$, so there are at most $C_L$
colours. If $j<j'$ have equal colours then
\[
 Z_\kappa=N(z^j_\kappa)-N(z^{j'}_\kappa)\ge0
\]
counts exactly original internal windows first touched during $(j,j']$.
Equal boundary tags give equal endpoint imbalance, so each $Z_\kappa$
is a circulation (identical short words give zero). Its nonzero support
is in the single $J$ just described. Crossing windows in $W_j,W_{j'}$
agree by the boundary identity. Thus $b(W_j)=b(W_{j'})$ and
$N(W_j)-N(W_{j'})=\sum_\kappa Z_\kappa$.

From each visited index $j$, go to the last $j'\ge j$ of its colour, delete
the first-touches in $(j,j']$, then, if present, retain the single increment
to $j'+1$ and continue there. Starting colours never repeat. The jump
windows are disjoint across jumps and across maximal free subtrees.
The prescription concerns original positions, not positions in a subsequently
rearranged word. A zero jump needs no substitution.

\emph{Numerical rerouting.}
In a strongly connected $J$ suppose $m$ designated paths have prescribed
respective endpoints $p_l,q_l$ and combined target counts $M$ with
$\partial M=\sum_l(\mathbf e_{q_l}-\mathbf e_{p_l})$.
To expose a closed walk $g_0$ at vertex $u$ inside a designated main path,
it suffices that $M(e)\ge(m-1)+1+N_{\rm edge}(g_0)(e)+1$ on $J$.
Choose a simple path in each non-main slot and from the main start to $u$.
Subtract them and the prescribed closed walk. The remainder has every
edge positive, and imbalance from $u$ to the main end, so admits an Euler
trail (add a return edge, splice cycles in the connected balanced graph,
then remove that edge). Use it after $g_0$. All intervening frozen edges
remain genuine windows because their entering and leaving vertices were
preserved. The Euler remainder is nonempty. Empty non-main paths are allowed.

\emph{One deletion.}
Maintain for each root component its original $b,c$ and current vector $M_i$
obtained by subtracting earlier jumps. Maintain that there is a realizing
tuple, and all such tuples are equivalent to $x$. This holds initially by
the context invariant. Group nonzero $Z_\kappa$ by original component and
frequent component $(i,J)$ to form $Z_{i,J}$. The number $e_0$ of active
pools is at most the arity of $\eta$, hence at most $f$.

In $W_{j'}$ permit only original frequent windows internal to $x_i$ with
all positions outside the component intervals of $\eta$. They survive as
consecutive original letters. Freeze all other edges, including boundary
windows in the sentence. Removing the windows meeting one interval of
letters removes an interval of window starts; this increases the number of
remaining runs of a given pool by at most one. Thus the number $m$ of
flexible runs is at most $R+1+f$. Every selected occurrence of a frequent
edge is wholly outside $\eta$ since every one of its letters is marked.
Thus at least $c_{g+1}$ occurrences of each such edge are flexible.
Marking alone does not make a rare edge flexible. For each active pool
choose the first original frequent interval containing a flexible window
as anchor; use the first flexible run there on the sentence side and the
whole corresponding interval on the representative side.

Fix a vertex $u$ of $J$ and a tour $D$ from $u$ covering all its edges of
length at most $\lambda$ (for each edge go to its tail, cross it, return
by simple paths). If $q>0$, put $D_0=D^q$, $h=D_0D_0$ and
$t=D_0\operatorname{Euler}_u(N_{\rm edge}(D_0)+Z_{i,J})$.
This Euler circuit exists by full support and balance. If $q=0$, use the
single vertex, take $h$ to traverse every loop of $J$ once and $t$ to
traverse $N_{\rm edge}(h)+Z_{i,J}$. Write $h,t$ also for the emitted
last-letter strings. They are positive and have the same content, namely
the set of emitted labels on $J$. For $q>0$ both are long, share their first
$q$ letters, and end in $u$. For $q=0$ the window type is trivial.
Their component $s$-types agree. Their edge-count difference is $Z_{i,J}$;
when comparing internal word counts the first $q$ emitted letters account
for identical omitted incoming edges because the initial vertex and those
letters agree.

On the representative side keep the original prefix vertex and rare edge
slots of each component. In every frequent pool the current restriction
has the total imbalance of its original intervals, since earlier deletions
were circulations. It contains, in addition to the present $Z_{i,J}$,
at least $c_{g+1}$ copies per edge: selected occurrences and present jump
occurrences are disjoint and neither has been deleted. Expose $t$ in the
anchor of every active pool, using numerical rerouting; inactive pools
use an empty prescribed walk. Call the tuple $y$. Replacing the exposed
$t$ by $h$ gives $y'$ with counts decreased by exactly the jump.
Content is preserved by the original prefix and the support of emitted
edges. End vertices persist, and a modified component remains long since
a frequent edge retains at least $c_{g+1}$ occurrences ($>2q$ for $q>0$).
Components with no frequent edges are kept verbatim.

On the sentence side use just the flexible runs in each active pool.
Their actual edge counts have the required imbalance. Expose $h$ at the
anchor, freeze the rest of the sentence outside these runs, obtaining
$F[h]$. In both applications the numerical condition follows, edgewise,
from $m\le R+1+f$ and
\[
 c_{g+1}-(R+1+f+ B+1)=c_g+K-1\quad(q>0).
\]
Here $N_{\rm edge}(h)(e)\le B$; on the representative side the additional
$Z_{i,J}(e)$ pays for $t$. For $q=0$, $B=0$ but the per-edge cost of
$h$ is one, and the slack $c_g+K-1\ge1$ pays this cost on both sides.
The same bound covers inactive pools with a zero prescribed walk.

The sentence $F[h]$ has the length, end windows and counts of $W_{j'}$;
$F[t]$ has those of $W_j$ by the colour identity and endpoint-preserving
replacement. In a nonzero jump the anchors ensure long sentences;
for $q=0$ positive suffices. Both sentences therefore belong to $L$.
The holes on both sides are in the order of root components and their
original anchor intervals. Within one root component different anchors
of different pools have a frozen rare edge between them (its emitted
letter survives). Between root components the outer context has a positive
separator. On the sentence side such a separating rare edge is intact
since all its original positions were marked. Alternatively every main
path has a nonempty remainder after its hole. Thus both $F$ and the
context $F_E$ obtained by cutting $t$ out of $E[y]$ are admissible for
every admissible $E$; outside gaps may be empty.
The comparison tuples are the ordered lists $t,h$ just constructed. The
shared context is $F$, the test context $F_E$; the four words are
$F[t],F[h],F_E[t]=E[y],F_E[h]=E[y']$. Substitutability gives equality
of the last two membership indicators for every $E$. Their root component
types agree; the context invariant extends the conclusion from $y'$ to
every realization of its new data. This proves the induction.

\emph{Length.}
Count surviving original occurrences numerically. A window not internal
to one free component either touches a skeleton-written terminal
(at most $\ell N_{\rm sk}\theta$ starts) or crosses an endpoint of a free
component (at most $2\ell frN_{\rm sk}$ starts). Each other surviving
window is assigned to its unique first-touch increment. At most $C_L$
increments per maximal free subtree were retained, each touching at most
$\ell H\theta$ windows. Every component has length at most $q$ plus its
window count (including the short case). The asserted bound follows.
Only the original grammar supplied the short substitutions.
\end{proof}

\paragraph{Two corollaries.}
For a strictly $\ell$-local skeleton $R$, prescribe allowed internal windows,
end conditions and the short-word cases. Let $g$ be a morphism to an abelian
group, $H_a$ a subgroup, and take $L=R\cap g^{-1}(\gamma+H_a)$.
If $L$ is an $f$-MCFL it satisfies the theorem for $s=(c,b)$.
Indeed take any arity up to $f$, equally typed positive tuples $x,y$,
a shared accepting context $E_0$, and any test context $E_1$. In the
four words $E_0[x],E_0[y],E_1[x],E_1[y]$ the internal windows of either
tuple are allowed by the first two words. For the last two, end data,
short-word cases, fixed-piece windows and crossing windows agree, so
skeleton membership agrees. The abelian value change is the same sum
of component differences in both contexts, hence lies in $H_a$.
The required local-count description follows since letter counts are read
from the initial vertex and emitted edge labels (or directly from short
words). For finite abelian groups the intersection is regular.

If $k\ge2$, choose different letters $a,b_0$. In the positive image of
$s=(c,b)$, $e=s(a^{2\ell})$ and $e s(b_0)e$ are distinct idempotents of
$e\,s(\Sigma^+)\,e$, as is seen from their contents and boundaries.
This local monoid is not a group.
\end{document}
